\chapter{Sviluppi futuri}
\label{cap:sviluppi-futuri}

Nel corso del lavoro sono state concepite alcune funzionalità che non sono state realizzate.
Non si tratta di omissioni né di lavoro rimasto incompiuto per mancanza di tempo: sono
\textbf{decisioni}, prese confrontando ciò che ciascuna avrebbe aggiunto con ciò che sarebbe
costata, e va reso conto tanto delle une quanto delle altre. Un progetto che documenta solo
ciò che ha costruito racconta metà del proprio percorso, e non consente a chi lo riprende di
capire se una strada sia stata esclusa perché senza sbocco o semplicemente non tentata. \medskip

Il criterio applicato non è stato «questa funzionalità è utile?», domanda a cui tutte e tre
rispondono affermativamente, bensì il concorso di tre condizioni: il \textbf{valore} che la
funzionalità porta all'utilizzatore, il \textbf{costo} della sua realizzazione, e il
\textbf{contributo metodologico}, ossia se essa produca un risultato misurabile e discutibile
in sede di tesi oppure sia lavoro di realizzazione che nulla aggiunge alla domanda di
ricerca. È soprattutto la terza condizione a spiegare gli esiti che seguono: a parità di
utilità pratica, una funzionalità che non genera alcun dato da valutare compete con il
tempo destinato a ciò che invece ne genera. \medskip

Le sezioni che seguono trattano una funzionalità ciascuna secondo lo stesso schema --- che
cosa avrebbe portato, che cosa sarebbe costata, perché è stata accantonata e a quali
condizioni potrebbe essere ripresa --- e la Tabella~\ref{tab:funzionalita-scartate} ne
riassume la valutazione. L'ultima sezione raccoglie invece i limiti già dichiarati altrove
nel documento, che restano aperti come naturale prosecuzione del lavoro svolto.

\section{Analisi contestuale della struttura del file system}
\label{sec:sf-analisi-struttura}
L'idea, formulata fra i requisiti funzionali (cfr.~\ref{sub:analisi-filesystem}) e prevista
in architettura come blocco autonomo (cfr.~\ref{sec:b2-non-realizzato}), era di trattare
l'\textbf{organizzazione delle cartelle} come un segnale: i nomi dei rami, la profondità e
il raggruppamento dei documenti dicono qualcosa sulla finalità per cui quei documenti sono
tenuti insieme. Nella forma più economica immaginata, l'ingresso sarebbe stato la resa
testuale dell'albero prodotta da un comune comando di sistema --- poche decine di righe di
testo per un'intera condivisione --- da sottoporre a un modello linguistico. \medskip

Il valore atteso era duplice: \textbf{proporre} automaticamente l'associazione fra rami
dell'albero e attività di trattamento, che oggi il DPO stabilisce a mano, e segnalare le
\emph{anomalie posizionali}, ossia i documenti il cui contenuto non è coerente con la
cartella che li ospita. È un valore reale, e resta l'unico modo per far emergere una
violazione del principio di limitazione della finalità che non sia già visibile nel
confronto fra dati rilevati e dati dichiarati. \medskip

Il costo, però, non sta dove sembra. Procurarsi l'ingresso è banale; interpretarlo non lo è.
Un'inferenza sul significato di una struttura di cartelle è un'euristica, e un'euristica va
\textbf{misurata}: servirebbe un corpus di alberi documentali annotati con l'attribuzione
«corretta» di ciascun ramo, che non esiste e andrebbe costruito sinteticamente, raddoppiando
di fatto l'apparato di valutazione (cfr.~\ref{sub:rischio-dataset}) per una funzionalità
accessoria. A ciò si aggiunge un problema di merito: un'attribuzione inferita dalla macchina
diventa il presupposto del verdetto di conformità, che verrebbe così a poggiare su una
supposizione anziché su una dichiarazione. \medskip

La ragione decisiva dell'accantonamento è tuttavia un'altra: la parte \emph{utile} della
funzionalità è stata ottenuta per altra via. Le regole che associano un prefisso di percorso
a una o più attività (cfr.~\ref{sub:b6-associazione}) producono lo stesso effetto operativo
--- associare un intero ramo in un solo gesto, inclusi i documenti che vi compariranno in
seguito --- con la differenza sostanziale che a stabilire il significato della collocazione è
il DPO. Ciò che resterebbe da realizzare non è quindi l'associazione, ma la sola
\emph{proposta} di quelle regole: un incremento reale ma marginale, che presuppone comunque
l'apparato di valutazione appena descritto. La funzionalità è perciò rinviata, e la sua forma
naturale di ripresa è quella di un proponente che suggerisca regole al DPO senza mai
applicarle da sé.

\section{Canale di segnalazione dai dipendenti}
\label{sec:sf-feedback-loop}
Il secondo requisito rimasto sulla carta (cfr.~\ref{sub:feedback-loop}) prevedeva un canale
attraverso cui il personale dell'organizzazione potesse segnalare al DPO i dati personali
incontrati nel lavoro quotidiano, sfuggiti al rilevamento automatico. Il DPO avrebbe
valutato le segnalazioni e le avrebbe usate per arricchire il registro dei trattamenti o i
parametri di rilevamento, secondo il ciclo illustrato nella
Figura~\ref{fig:human-in-the-loop}.
\begin{figure}[htbp]
	\centering
	\begin{tikzpicture}[
			sysnode/.style={rectangle, draw, rounded corners, minimum width=3.6cm,
					minimum height=1.1cm, align=center, fill=blue!8, font=\small},
			human/.style={rectangle, draw, rounded corners, minimum width=3.4cm,
					minimum height=1.1cm, align=center, fill=purple!8, font=\small},
			store/.style={cylinder, draw, shape border rotate=90, aspect=0.25,
					minimum height=1.5cm, minimum width=2.6cm, align=center, fill=gray!10, font=\small},
			param/.style={rectangle, draw, rounded corners, minimum width=3.6cm,
					minimum height=1.1cm, align=center, fill=orange!10, font=\small},
			io/.style={rectangle, draw, rounded corners, minimum width=3.2cm,
					minimum height=0.9cm, align=center, fill=green!8, font=\small},
			arrow/.style={thick, -{Latex[length=2mm]}},
		]

		% --- Nodi ---
		% I documenti sono l'artefatto condiviso: analizzati dal sistema e, in
		% parallelo, consultati dai dipendenti nel loro lavoro ordinario. I
		% dipendenti restano ignari del rilevamento: per loro sono solo documenti.
		\node (docs)  [io]      at (3,    1.4)  {Documenti aziendali};
		\node (sys)   [sysnode] at (-0.5, -1.4) {Sistema di rilevamento\\PII};
		\node (emp)   [human]   at (6.5,  -1.4) {Dipendenti\\dell'azienda};
		\node (dpo)   [human]   at (6.5,  -4.2) {DPO\\(analisi segnalazioni)};
		\node (ropa)  [store]   at (-0.5, -3.8) {ROPA\\(trattamenti,\\categorie)};
		\node (param) [param]   at (-0.5, -6.0) {Parametri del sistema\\(catalogo PII, regole)};

		% --- I documenti: analizzati dal sistema, consultati dai dipendenti ---
		\draw [arrow] (docs) -| node[pos=0.25, above, font=\footnotesize]{analizza} (sys.north);
		\draw [arrow] (docs) -| node[pos=0.25, above, font=\footnotesize, align=center]{consultati nel\\lavoro quotidiano} (emp.north);

		% --- Segnalazione: dipendenti -> DPO ---
		\draw [arrow] (emp) -- node[right, font=\footnotesize, align=left]{segnalano le PII\\in cui si imbattono} (dpo);

		% --- Fork: il DPO valuta e arricchisce ROPA e parametri ---
		\coordinate (fk) at (3, -4.2);
		\draw [thick] (dpo.west) -- node[above, font=\footnotesize]{valuta e arricchisce} (fk);
		\draw [arrow] (fk) |- (ropa.east);
		\draw [arrow] (fk) |- (param.east);

		% --- Gli aggiornamenti rientrano nel sistema, che analizza meglio i documenti ---
		\draw [arrow] (ropa) -- (sys);
		\draw [arrow] (param.west) -- ++(-1.3,0) |- (sys.west);

	\end{tikzpicture}
	\caption{Meccanismo di \emph{human-in-the-loop}: nel normale lavoro con i
		documenti aziendali i dipendenti segnalano al DPO le PII in cui si imbattono;
		il DPO analizza le segnalazioni e arricchisce il ROPA o i parametri del
		sistema, migliorando il rilevamento. Il processo di analisi resta trasparente
		ai dipendenti, per i quali i file sono semplici documenti.}
	\label{fig:human-in-the-loop}
\end{figure}
Il valore atteso era il più alto dei tre, e mirava al punto più debole del sistema. I falsi
negativi sono per definizione invisibili (cfr.~\ref{sub:rischio-falsi-negativi}): nessuna
misura interna può rivelare un dato personale che il rilevatore non vede. I dipendenti sono
l'unico sensore che osserva i documenti nel loro contesto d'uso, e sono quindi l'unica
sorgente capace di chiudere quella cecità dall'esterno. \medskip

Il costo è però sproporzionato, e per una ragione strutturale più che di quantità di lavoro.
Il sistema è progettato per un \textbf{unico profilo di utilizzatore}, il DPO, e su questa
premessa poggia l'intera restituzione: una sola interfaccia, che espone senza distinzioni
tutto ciò che il sistema sa. Occorre qui evitare un equivoco: proteggere quella console con
un'autenticazione è cosa dovuta e di costo trascurabile, ed è annoverata fra gli interventi
da compiere (cfr.~\ref{sec:sf-direzioni}); ciò che è oneroso non è \emph{autenticare}, ma
\textbf{autorizzare}. Aprire un canale al personale significherebbe passare da un'identità a
un insieme di identità da censire e mantenere allineato all'organico, e con esso introdurre
ruoli, permessi differenziati, un'interfaccia distinta per chi non deve vedere la console, la
notifica e la moderazione delle segnalazioni. È questa la parte preponderante del lavoro,
mentre la funzionalità in sé si riduce a un modulo di inserimento. \medskip

Vi è inoltre un argomento che da solo sconsiglia la realizzazione nella forma prevista. Il
registro delle rilevazioni non contiene i valori dei dati personali, ma ne costituisce la
\textbf{mappa}: dice dove si trovano, ed è per questo un artefatto sensibile
(cfr.~\ref{sub:rischio-sicurezza-dati}). Esporre a tutto il personale un'interfaccia
affacciata su quel registro significherebbe ampliarne la superficie di accesso dalla singola
figura del DPO all'intera organizzazione: un peggioramento del livello di sicurezza in uno
strumento che esiste per ridurre l'esposizione dei dati. \medskip

A questo si aggiunge la considerazione che ha determinato l'esito: la funzionalità è un
\textbf{processo organizzativo}, non una tecnica di rilevamento. Non produce alcun risultato
misurabile --- non migliora una metrica, non consente un confronto, non sostiene né confuta
un'ipotesi --- e il suo effetto dipende interamente dalla partecipazione delle persone, che
un prototipo non può mettere alla prova. È lavoro di realizzazione che non alimenta la
domanda di ricerca. La funzionalità è pertanto esclusa dal perimetro, e non semplicemente
rinviata; qualora venisse ripresa, la forma corretta sarebbe un canale di sola immissione,
separato dalla console del DPO e privo di qualsiasi accesso in lettura al registro.

\section{Estensione dei formati documentali}
\label{sec:sf-formati}
Il requisito di estrazione (cfr.~\ref{sub:estrazione-testo}) è soddisfatto sui formati
\textit{born-digital} dell'ufficio --- documenti impaginati, videoscrittura, fogli di calcolo
e testo semplice. Restano fuori i documenti che sono immagini, ossia le scansioni, e con essi
il riconoscimento ottico dei caratteri, i layout complessi a più colonne e le tabelle, oltre
ai formati meno diffusi. Allo stato attuale se ne fa a meno: i lettori impiegati sono
sufficienti ai formati coperti, e nessun motore di analisi documentale più elaborato è stato
introdotto. Qualora l'insieme dei formati da trattare venisse ampliato, sarebbe questo il
punto in cui valutare l'adozione di uno strumento dedicato --- Docling o analoghi --- in
luogo dei lettori specifici per singolo formato. \medskip

Il valore è reale ma \textbf{lineare}: ogni formato aggiunto rende leggibile un'altra fetta
del patrimonio documentale, e nulla più. Non cambia ciò che il sistema sa fare, ma solo su
quanti documenti riesce a farlo --- il che, va detto, in un archivio reale con anni di
scansioni può non essere poco. \medskip

Anche il costo è in larga parte lineare, ed è per questo che la funzionalità è stata giudicata
senza contributo di ricerca: lavoro meccanico, ripetitivo e senza decisioni
interessanti da prendere. Il riconoscimento ottico fa però eccezione, perché introduce un
salto qualitativo: non è un lettore in più, è una \textbf{nuova sorgente di errore} collocata
a monte del rilevamento. Il testo che ne esce contiene errori di trascrizione, e ogni PII non
rilevata diventerebbe attribuibile indifferentemente al rilevatore o all'estrattore. \medskip

È quest'ultimo il motivo per cui l'estensione non è stata affrontata ora: la valutazione
misura deliberatamente la \emph{tassa di estrazione} come grandezza distinta dalle prestazioni
del rilevamento (cfr.~Capitolo~\ref{cap:test}), e introdurre un estrattore rumoroso avrebbe
confuso le due sorgenti di errore proprio mentre le si stava separando. La funzionalità resta
rinviata ed è, delle tre, quella ripresa con il minor rischio: il perimetro è contenuto, il
contratto di estrazione già in essere non cambia, e l'ordine ragionevole è affrontare prima i
formati testuali residui e solo dopo, isolandone la misura, il riconoscimento ottico.

\section{Quadro riassuntivo}
\label{sec:sf-quadro}

\begin{table}[htbp]
	\centering
	\caption{Valutazione delle funzionalità concepite e non realizzate.}
	\label{tab:funzionalita-scartate}
	\small
	\resizebox{\textwidth}{!}{%
		\begin{tabular}{|l|c|c|l|c|}
			\hline
			\textbf{Funzionalità}                        & \textbf{Valore} & \textbf{Costo} & \textbf{Contributo metodologico}     & \textbf{Esito}   \\
			\hline
			Analisi della struttura del file system (B2) & Medio           & Alto           & Basso: la parte utile è già ottenuta & Rinviata         \\
			\hline
			Segnalazione dai dipendenti                  & Alto            & Alto           & Nullo: processo organizzativo        & \textbf{Esclusa} \\
			\hline
			Formati ulteriori e riconoscimento ottico    & Medio           & Medio          & Negativo: confonde la misura         & Rinviata         \\
			\hline
		\end{tabular}%
	}
\end{table}

Il quadro mette in evidenza un fatto che le singole sezioni lasciano implicito: nessuna delle
tre è stata scartata perché di poco valore. La segnalazione dai dipendenti è anzi quella dal
valore più alto, ed è l'unica esclusa in via definitiva. A determinare gli esiti è stata la
terza colonna: dove la funzionalità non produce un risultato valutabile, il suo costo compete
direttamente con il tempo destinato a ciò che invece lo produce.

\section{Limiti dichiarati e direzioni residue}
\label{sec:sf-direzioni}
Oltre alle funzionalità sopra discusse, restano aperti alcuni limiti già riconosciuti nel
corso del documento, che costituiscono la naturale prosecuzione del lavoro:

\begin{itemize}
	\item \textbf{Autenticazione della console.} L'interfaccia del DPO è oggi raggiungibile da
	      chiunque abbia accesso di rete alla porta su cui è servita. Anche per un applicativo a
	      utilizzatore singolo, e per quanto collocato in una rete interna, l'accesso va protetto
	      da una schermata di autenticazione: è prassi di base, e a maggior ragione qui, dove ciò
	      che si protegge è la mappa della collocazione dei dati personali
	      (cfr.~\ref{sub:rischio-sicurezza-dati}). L'intervento è circoscritto e di costo
	      trascurabile, e va tenuto distinto dalla gestione di più profili discussa in
	      \ref{sec:sf-feedback-loop}: qui si tratta di verificare \emph{un'unica} identità, non di
	      amministrarne molte.
	\item \textbf{Parallelizzazione dell'analisi.} L'elaborazione incrementale riduce il costo
	      delle esecuzioni successive alla prima, ma non quello della passata iniziale, che resta
	      sequenziale (cfr.~\ref{sub:rischio-scalabilita}).
	\item \textbf{Validazione su corpus reale.} Le misure disponibili provengono da corpus
	      sintetici, più ordinati del reale, e vanno intese come limite superiore; la verifica sul
	      patrimonio documentale effettivo di un'organizzazione resta un passaggio non eludibile
	      prima dell'adozione (cfr.~\ref{sub:rischio-dataset}, \ref{sub:rischio-domain-silo}).
	\item \textbf{Passata generativa selettiva.} Il punto di innesto per un rilevatore basato su
	      modello generativo, da impiegare su un campione di documenti, è previsto e predisposto
	      ma non ancora attivo (cfr.~\ref{sub:b4-ai-previsto}).
	\item \textbf{Datazione dal contenuto.} La data di riferimento di un documento è oggi stimata
	      dai metadati e dal file system; ricavarla dal corpo del testo --- dove spesso compare in
	      chiaro --- renderebbe la verifica della conservazione assai più solida, ma richiede un
	      riconoscimento delle date e una disambiguazione fra le molte presenti in un documento.
\end{itemize}
